% uid: ""
% school: "lycee"
% level: "3eme"
% section: ["Economie et Gestion"]
% chapter: "Suites réelles"
% title: "Dons décroissants d'une association"
% difficulty: "Moyen"
% points: null
% tags: ["suites", "somme-geometrique", "limite", "seuil", "modelisation"]
% provenance: "own"
% dressing: "somme cumulée"
% rung: 3
% batch: "2026-11-20-suites-3eco-A"

% ---STATEMENT---
Une association reçoit d'un mécène un don de $600$ dinars la première année ($D_0=600$). Chaque année suivante, le don est égal à $75\,\%$ du don de l'année précédente. On note $D_n$ le don de l'année $n$ et $C_n=D_0+D_1+\dots+D_n$ le total reçu.

\begin{enumerate}
  \item Exprimer $D_n$ en fonction de $n$.
  \item Montrer que $C_n=2\,400\,\bigl(1-0{,}75^{\,n+1}\bigr)$.
  \item Déterminer $\displaystyle\lim_{n\to+\infty}C_n$. L'association peut-elle espérer recevoir $2\,500$ dinars au total ?
  \item À partir de quelle année le total dépasse-t-il $2\,300$ dinars ?
\end{enumerate}

% ---SOLUTION---
\begin{enumerate}
  \item Chaque don est $75\,\%$ du précédent : $D_{n+1}=0{,}75\,D_n$. La suite $(D_n)$ est géométrique de raison $0{,}75$ et de premier terme $600$ :
  $$ \ib{$D_n=600\times0{,}75^n$}.$$
  \item $C_n=D_0+D_1+\dots+D_n$ est la somme de $n+1$ termes d'une suite géométrique de raison $q=0{,}75\neq1$ :
  \begin{prop}[Somme géométrique]
  $u_0+\dots+u_n=u_0\,\dfrac{1-q^{\,n+1}}{1-q}$.
  \end{prop}
  $$C_n=600\times\frac{1-0{,}75^{\,n+1}}{1-0{,}75}=\frac{600}{0{,}25}\bigl(1-0{,}75^{\,n+1}\bigr)=\ib{$2\,400\bigl(1-0{,}75^{\,n+1}\bigr)$}.$$
  \item $0<0{,}75<1$, donc $0{,}75^{\,n+1}\to0$ et $C_n\to2\,400\times(1-0)=$ \ib{$2\,400$}.

  Comme $C_n=2\,400\bigl(1-0{,}75^{\,n+1}\bigr)<2\,400$ pour tout $n$, le total ne dépassera \emph{jamais} $2\,400$ dinars : l'association ne peut pas espérer $2\,500$ dinars.
  \item Le total est croissant (on ajoute des dons positifs). On calcule :
  $$C_{10}=2\,400\bigl(1-0{,}75^{11}\bigr)\approx2\,298{,}6\qquad C_{11}=2\,400\bigl(1-0{,}75^{12}\bigr)\approx2\,324{,}0.$$
  $C_{10}<2\,300<C_{11}$ : le total dépasse $2\,300$ dinars à partir de l'année $n=11$.

  \textbf{Remarque :} algébriquement, $C_n>2\,300\iff0{,}75^{\,n+1}<\dfrac1{96}$.
\end{enumerate}

\textbf{Idée :} la somme cumulée d'une suite géométrique de raison $q\in]0\,;1[$ reste inférieure à $\dfrac{u_0}{1-q}$.
